%! Inverse Functions — Unit 5, Lesson 5.7 — Slides (Beamer)
\documentclass[aspectratio=169, 11pt]{beamer}
\usepackage{algebra2-beamer}   % provides \forestheader{} and \sectionlabel[color]{}
% Standard coordinate grid + axes: {xmin}{xmax}{ymin}{ymax}
\newcommand{\lgrid}[4]{%
  \draw[step=1, linegray!45, thin] (#1,#3) grid (#2,#4);
  \draw[->, thick, charcoal] (#1,0)--(#2,0) node[below right, font=\tiny]{$x$};
  \draw[->, thick, charcoal] (0,#3)--(0,#4) node[left, font=\tiny]{$y$};}

\begin{document}

% TITLE SLIDE (CourseName is not defined in beamer, so the name is literal)
{
\setbeamercolor{background canvas}{bg=forest}
\begin{frame}[plain]
  \vspace{2.8cm}\hspace{0.55cm}%
  \begin{minipage}{0.9\textwidth}
    {\color{goldacc}\bfseries\small LESSON 5.7}\\[0.5cm]
    {\color{white}\bfseries\LARGE Inverse Functions}\\[1.2cm]
    {\color{forestmid}\small Algebra 2: Shepherd~$\cdot$~Unit 5: Radical Functions}
  \end{minipage}
\end{frame}
}

% HOOK
\begin{frame}[t, plain]
  \forestheader{Two students, one formula, two answers}
  \sectionlabel{Hook}
  \begin{columns}[T]
    \begin{column}{0.44\textwidth}
      \footnotesize
      Celsius to Fahrenheit:
      \[ F(c)=\tfrac{9}{5}c+32 \]
      \emph{multiply by $\tfrac95$, then add $32$.}
      \vspace{4pt}
      \begin{center}
      \renewcommand{\arraystretch}{1.3}
      \begin{tabular}{|c|c|c|c|}
      \hline
      $c$ & $0$ & $100$ & $20$ \\ \hline
      $F(c)$ & $32$ & $212$ & $68$ \\ \hline
      \end{tabular}
      \end{center}
      \vspace{6pt}
      Now go the \emph{other} way.
    \end{column}
    \begin{column}{0.52\textwidth}
      \footnotesize
      \begin{center}
      \renewcommand{\arraystretch}{1.5}
      \begin{tabular}{|l|l|c|}
      \hline
      \textbf{Rosa} & $C(f)=\frac{5}{9}f-32$ & at $212$: $\approx85.8$ \\ \hline
      \textbf{Malik} & $C(f)=\frac{5}{9}(f-32)$ & at $212$: $\mathbf{100}$ \\ \hline
      \end{tabular}
      \end{center}
      \vspace{6pt}
      \begin{block}{Rosa reversed both operations correctly}
        What she did \emph{not} reverse was the \textbf{order}. \\[3pt]
        Same mistake as: undoing ``double, then add $5$'' by halving first.
      \end{block}
    \end{column}
  \end{columns}
  \vspace{3pt}
  \begin{center}
    \textbf{To undo a process, reverse the operations \emph{and} the order. Shoes before socks.}
  \end{center}
\end{frame}

% THE NOTATION
\begin{frame}[t, plain]
  \forestheader{A function that undoes another}
  \sectionlabel[goldacc]{The notation}
  \begin{center}
  \begin{tikzpicture}[every node/.style={font=\small}]
    \node (x) at (0,0) {$x$};
    \node[draw=forest, very thick, rounded corners, fill=forest!12, minimum width=1.5cm,
          minimum height=0.9cm] (f) at (3.0,0) {$f$};
    \node[draw=navy, very thick, rounded corners, fill=navy!12, minimum width=1.5cm,
          minimum height=0.9cm] (fi) at (7.4,0) {$f^{-1}$};
    \node (out) at (10.4,0) {$x$};
    \draw[->, very thick, charcoal] (x)--(f);
    \draw[->, very thick, charcoal] (f)--node[above, font=\scriptsize]{$f(x)$}(fi);
    \draw[->, very thick, charcoal] (fi)--(out);
  \end{tikzpicture}
  \end{center}
  \vspace{4pt}
  \begin{columns}[T]
    \begin{column}{0.48\textwidth}
      \footnotesize
      \begin{block}{The $-1$ is \emph{not} an exponent}
        $f^{-1}(x)$ is a \textbf{label}: ``the undoing function.'' \\[4pt]
        $f^{-1}(x) \ne \dfrac{1}{f(x)}$ \quad --- ever.
      \end{block}
    \end{column}
    \begin{column}{0.48\textwidth}
      \footnotesize
      \begin{center}
      \renewcommand{\arraystretch}{1.5}
      \begin{tabular}{|l|c|c|}
      \hline
      $f(x)=2x$ & rule & at $x=4$ \\ \hline
      $f^{-1}(x)$ & $\frac{x}{2}$ & $2$ \\ \hline
      $\frac{1}{f(x)}$ & $\frac{1}{2x}$ & $\frac18$ \\ \hline
      \end{tabular}
      \end{center}
      \vspace{4pt}
      \centering Two completely different functions.
    \end{column}
  \end{columns}
\end{frame}

% THE TEST
\begin{frame}[t, plain]
  \forestheader{The test: compose \emph{both} ways and look for $x$}
  \sectionlabel{Justify it}
  \begin{center}
    \textbf{$f$ and $g$ are inverses} \quad exactly when \quad
    $f\big(g(x)\big)=x$ \; \textbf{and} \; $g\big(f(x)\big)=x$.
  \end{center}
  \vspace{6pt}
  \begin{columns}[T]
    \begin{column}{0.48\textwidth}
      \footnotesize
      \textbf{Passes.} $f(x)=2x+6$, \; $g(x)=\frac{x-6}{2}$
      \[ 2\!\left(\tfrac{x-6}{2}\right)+6=x \qquad \tfrac{(2x+6)-6}{2}=x \]
      \vspace{2pt}
      \textbf{Fails.} $h(x)=\frac{x}{2}-6$ --- right operations, wrong order
      \[ f\big(h(x)\big)=2\!\left(\tfrac{x}{2}-6\right)+6=x-6 \]
      \centering \emph{Off by a constant --- only composition catches it.}
    \end{column}
    \begin{column}{0.48\textwidth}
      \footnotesize
      \begin{block}{Why \emph{both} orders? Yesterday's cliffhanger}
        $f(x)=x^{2}$, \; $g(x)=\sqrt{x}$: \\[4pt]
        $\left(\sqrt{x}\,\right)^{2}=x$ \quad --- but only for $x\ge0$ \\[3pt]
        $\sqrt{x^{2}}=|x|$ \quad --- which is \emph{not} $x$ when $x<0$ \\[4pt]
        One order passing is \textbf{not} a verdict.
      \end{block}
    \end{column}
  \end{columns}
\end{frame}

% SWAP AND SOLVE
\begin{frame}[t, plain]
  \forestheader{Building the partner: swap and solve}
  \sectionlabel[goldacc]{The construction}
  \begin{center}
    \textbf{1.} Write $y=f(x)$. \quad \textbf{2.} Trade $x$ and $y$. \quad
    \textbf{3.} Solve for $y$. \quad \textbf{4.} Rename it $f^{-1}(x)$.
  \end{center}
  \vspace{4pt}
  \begin{columns}[T]
    \begin{column}{0.54\textwidth}
      \footnotesize
      \begin{center}
      \renewcommand{\arraystretch}{1.45}
      \begin{tabular}{|l|l|l|}
      \hline
      $f(x)$ & \textbf{swap} & $f^{-1}(x)$ \\ \hline
      $3x-7$ & $x=3y-7$ & $\frac{x+7}{3}$ \\ \hline
      $\frac{x}{4}+2$ & $x=\frac{y}{4}+2$ & $4x-8$ \\ \hline
      $x^{3}+1$ & $x=y^{3}+1$ & $\sqrt[3]{x-1}$ \\ \hline
      $\sqrt{x-2}$ & $x=\sqrt{y-2}$ & $x^{2}+2$, \; $x\ge0$ \\ \hline
      \end{tabular}
      \end{center}
    \end{column}
    \begin{column}{0.42\textwidth}
      \footnotesize
      Step 3 is nothing new --- it is \emph{solve for the other variable}, from the Warm-Up.
      \vspace{5pt}
      \begin{block}{The last row is the hard one}
        $x^{2}+2$ \textbf{without} $x\ge0$ is a different function --- and it fails the test.
        \\[3pt]
        $f$ never sent anything to $-3$, so $f^{-1}(-3)$ has no right to answer.
      \end{block}
    \end{column}
  \end{columns}
\end{frame}

% REFLECTION
\begin{frame}[t, plain]
  \forestheader{Swap the equation $=$ swap every point}
  \sectionlabel{Reflection over $y=x$}
  \begin{columns}[T]
    \begin{column}{0.40\textwidth}
      \footnotesize
      If $(a,b)$ is on $f$, then $(b,a)$ is on $f^{-1}$.
      \vspace{5pt}
      \begin{center}
      \renewcommand{\arraystretch}{1.25}
      \begin{tabular}{|c|c|c|c|c|}
      \hline
      $x$ & $0$ & $1$ & $4$ & $9$ \\ \hline
      $\sqrt{x}$ & $0$ & $1$ & $2$ & $3$ \\ \hline
      \end{tabular}
      \\[6pt]
      \renewcommand{\arraystretch}{1.25}
      \begin{tabular}{|c|c|c|c|c|}
      \hline
      $x$ & $0$ & $1$ & $2$ & $3$ \\ \hline
      $f^{-1}(x)$ & $0$ & $1$ & $4$ & $9$ \\ \hline
      \end{tabular}
      \end{center}
      \vspace{5pt}
      The rows just traded places. \\
      \textbf{A2.F.2j:} $f^{-1}$ is the mirror image of $f$ across $y=x$.
    \end{column}
    \begin{column}{0.56\textwidth}
      \centering
      \begin{tikzpicture}[scale=0.52]
        \lgrid{-2}{6}{-2}{6}
        \draw[dashed, very thick, charcoal] (-1.6,-1.6)--(5.6,5.6);
        \begin{scope}
          \clip (-2,-2) rectangle (6,6);
          \draw[very thick, forest, domain=0:6, samples=140, smooth] plot (\x,{sqrt(\x)});
          \draw[very thick, navy, domain=0:2.5, samples=140, smooth] plot (\x,{(\x)*(\x)});
        \end{scope}
        \fill[forest] (0,0) circle (3.4pt);
        \node[forest, font=\scriptsize] at (5.0,1.5) {$y=\sqrt{x}$};
        \node[navy, font=\scriptsize] at (3.8,5.0) {$y=x^{2}$};
        \node[charcoal, font=\scriptsize] at (5.0,3.6) {$y=x$};
      \end{tikzpicture}
    \end{column}
  \end{columns}
\end{frame}

% POINTS ON y = x
\begin{frame}[t, plain]
  \forestheader{Intercepts swap. Points on $y=x$ never move.}
  \sectionlabel[goldacc]{Reading the picture}
  \begin{columns}[T]
    \begin{column}{0.50\textwidth}
      \centering
      \begin{tikzpicture}[scale=0.40]
        \lgrid{-6}{6}{-6}{6}
        \draw[dashed, very thick, charcoal] (-5.6,-5.6)--(5.6,5.6);
        \begin{scope}
          \clip (-6,-6) rectangle (6,6);
          \draw[very thick, forest] (-1,-6)--(6,8);
          \draw[very thick, navy] (-6,-1)--(8,6);
        \end{scope}
        \fill[charcoal] (4,4) circle (4pt);
        \node[forest, font=\scriptsize] at (-3.4,-2.2) {$f(x)=2x-4$};
        \node[navy, font=\scriptsize] at (-3.4,1.8) {$f^{-1}$};
      \end{tikzpicture}
    \end{column}
    \begin{column}{0.46\textwidth}
      \footnotesize
      \begin{center}
      \renewcommand{\arraystretch}{1.4}
      \begin{tabular}{|l|l|}
      \hline
      \textbf{on $f$} & \textbf{on $f^{-1}$} \\ \hline
      $(0,-4)$ & $(-4,0)$ \\ \hline
      $(2,0)$ & $(0,2)$ \\ \hline
      $(4,4)$ & $(4,4)$ \\ \hline
      \end{tabular}
      \end{center}
      \vspace{5pt}
      \begin{block}{Two consequences}
        The $x$-intercept of one is the $y$-intercept of the other. \\[3pt]
        $f$ and $f^{-1}$ can only meet \textbf{on the mirror line} --- which is why solving
        $f(x)=x$ finds it.
      \end{block}
    \end{column}
  \end{columns}
\end{frame}

% HORIZONTAL LINE TEST
\begin{frame}[t, plain]
  \forestheader{Does the inverse get to be a \emph{function}?}
  \sectionlabel{Horizontal line test}
  \begin{center}
    \begin{minipage}[c]{0.30\textwidth}
    \centering
    \begin{tikzpicture}[scale=0.34]
      \lgrid{-4}{4}{-5}{7}
      \begin{scope}
        \clip (-4,-5) rectangle (4,7);
        \draw[very thick, navy, domain=-2.7:2.7, samples=120, smooth] plot (\x,{(\x)*(\x)});
      \end{scope}
      \draw[very thick, redacc, dashed] (-4,4)--(4,4);
      \fill[redacc] (-2,4) circle (4.5pt);
      \fill[redacc] (2,4) circle (4.5pt);
    \end{tikzpicture}
    \\[2pt] {\scriptsize $y=x^{2}$ \\ \textbf{Fails} --- two hits}
    \end{minipage}
    \hfill
    \begin{minipage}[c]{0.30\textwidth}
    \centering
    \begin{tikzpicture}[scale=0.34]
      \lgrid{-4}{4}{-5}{7}
      \begin{scope}
        \clip (-4,-5) rectangle (4,7);
        \draw[very thick, navy, domain=0:2.7, samples=120, smooth] plot (\x,{(\x)*(\x)});
      \end{scope}
      \draw[very thick, greenacc, dashed] (-4,4)--(4,4);
      \fill[greenacc] (2,4) circle (4.5pt);
      \fill[navy] (0,0) circle (3.6pt);
    \end{tikzpicture}
    \\[2pt] {\scriptsize $y=x^{2},\ x\ge0$ \\ \textbf{Passes} --- restricted}
    \end{minipage}
    \hfill
    \begin{minipage}[c]{0.30\textwidth}
    \centering
    \begin{tikzpicture}[scale=0.34]
      \lgrid{-4}{4}{-5}{7}
      \begin{scope}
        \clip (-4,-5) rectangle (4,7);
        \draw[very thick, forest, domain=-1.75:1.9, samples=140, smooth]
          plot (\x,{(\x)*(\x)*(\x)});
      \end{scope}
      \draw[very thick, greenacc, dashed] (-4,4)--(4,4);
      \fill[greenacc] (1.587,4) circle (4.5pt);
    \end{tikzpicture}
    \\[2pt] {\scriptsize $y=x^{3}$ \\ \textbf{Passes} --- untouched}
    \end{minipage}
  \end{center}
  \vspace{4pt}
  \begin{center}
    $(-2,4)$ and $(2,4)$ swap into $(4,-2)$ and $(4,2)$ --- one input, two outputs. \\[3pt]
    \textbf{An inverse function exists exactly when no horizontal line crosses twice.}
  \end{center}
\end{frame}

% THE CAPSTONE
\begin{frame}[t, plain]
  \forestheader{The question Lesson 5.0 opened --- answered}
  \sectionlabel[goldacc]{The capstone}
  \begin{center}
  \footnotesize
  \renewcommand{\arraystretch}{1.4}
  \begin{tabular}{lll}
  \hline
   & \textbf{$y=x^{2}$ \; and \; $y=\sqrt{x}$} & \textbf{$y=x^{3}$ \; and \; $y=\sqrt[3]{x}$} \\
  \hline
  Arms & \textbf{two} & \textbf{one} \\
  Repeats an output? & yes ($-2$ and $2$ both give $4$) & never \\
  Horizontal line test & fails & passes \\
  Restriction needed? & \textbf{yes}, $x\ge0$ & \textbf{no} \\
  So the root's domain is & only $x\ge0$ & all real numbers \\
  Endpoint & yes, at the origin & none \\
  Extrema & absolute minimum & none \\
  \hline
  \end{tabular}
  \end{center}
  \vspace{8pt}
  \begin{center}
    \textbf{One sentence generates the whole table:} \\[3pt]
    $x^{2}$ has two arms and $x^{3}$ has one --- so the square must be cut down before it is allowed
    a partner, and the cube never has to be.
  \end{center}
\end{frame}

% DOMAIN AND RANGE
\begin{frame}[t, plain]
  \forestheader{Domain and range trade places}
  \sectionlabel{Where restrictions come from}
  \begin{columns}[T]
    \begin{column}{0.50\textwidth}
      \footnotesize
      \begin{center}
      domain of $f^{-1}$ $=$ \textbf{range} of $f$ \\[3pt]
      range of $f^{-1}$ $=$ \textbf{domain} of $f$
      \end{center}
      \vspace{6pt}
      \begin{center}
      \renewcommand{\arraystretch}{1.5}
      \begin{tabular}{|l|c|c|}
      \hline
       & \textbf{Domain} & \textbf{Range} \\ \hline
      $f(x)=\sqrt{x-2}$ & $x\ge2$ & $y\ge0$ \\ \hline
      $f^{-1}(x)=x^{2}+2$ & $x\ge0$ & $y\ge2$ \\ \hline
      \end{tabular}
      \end{center}
    \end{column}
    \begin{column}{0.46\textwidth}
      \footnotesize
      \begin{block}{The restriction is never guessed}
        It is the \textbf{range of $f$}, written down. \\[4pt]
        $\sqrt{x}-3$ has range $y\ge-3$ --- so its inverse takes $x\ge-3$, \emph{not} $x\ge0$.
      \end{block}
      \vspace{4pt}
      A student who writes $x\ge0$ every time has pattern-matched. The question is always:
      \emph{what does $f$ actually put out?}
    \end{column}
  \end{columns}
\end{frame}

% ERRORS
\begin{frame}[t, plain]
  \forestheader{Three ways to get it wrong --- and one number that catches each}
  \sectionlabel[goldacc]{Watch out}
  \begin{center}
  \footnotesize
  \renewcommand{\arraystretch}{1.5}
  \begin{tabular}{|l|l|l|}
  \hline
  \textbf{The error} & \textbf{Looks like} & \textbf{Caught by} \\ \hline
  Reversed operations, \emph{not} the order &
    $f(x)=3x+12 \Rightarrow \frac{x}{3}-12$ & $f(1)=15$; the inverse must return $1$ \\ \hline
  $-1$ read as an exponent &
    $f(x)=5x \Rightarrow \frac{1}{5x}$ & $f(3)=15$; that rule gives $\frac{1}{75}$ \\ \hline
  Restriction dropped &
    $\sqrt{x-1}\Rightarrow x^{2}+1$, no condition & test at $x=-2$: $f$ never output $-2$ \\ \hline
  \end{tabular}
  \end{center}
  \vspace{8pt}
  \begin{center}
    All three survive an eyeball check. \\[3pt]
    \textbf{None of them survives composing both ways.}
  \end{center}
\end{frame}

% SUMMARY
\begin{frame}[t, plain]
  \forestheader{Inverse functions, in one table}
  \sectionlabel{Summary}
  \begin{center}
  \footnotesize
  \renewcommand{\arraystretch}{1.35}
  \begin{tabular}{ll}
  \hline
  What it means & $f^{-1}$ undoes $f$: if $f$ sends $a$ to $b$, $f^{-1}$ sends $b$ to $a$ \\
  The notation & a label, never an exponent --- $f^{-1}(x)\ne\frac{1}{f(x)}$ \\
  How to prove it & compose \textbf{both} ways; both must give $x$ \\
  How to build it & swap $x$ and $y$, solve for $y$ --- reverse the operations \emph{and} the order \\
  What it looks like & the reflection of $f$ over the line $y=x$ \\
  Domain \& range & they trade places --- and that is where every restriction comes from \\
  When it exists & exactly when $f$ passes the horizontal line test \\
  \hline
  \end{tabular}
  \end{center}
  \vspace{8pt}
  \begin{center}
    $\left(\sqrt{x}\,\right)^{2}=x$ but $\sqrt{x^{2}}=|x|$. \\[3pt]
    \textbf{That is why the test needs both orders --- and why $y=x^{2}$ needs a restricted domain.}
  \end{center}
\end{frame}

% ROADMAP
\begin{frame}[t, plain]
  \forestheader{That is Unit 5 --- the loop is closed}
  \sectionlabel{Connections}
  \begin{columns}[T]
    \begin{column}{0.52\textwidth}
      \footnotesize
      \begin{itemize}\setlength{\itemsep}{3pt}
        \item \textbf{5.0} --- why does $\sqrt{x}$ get only half the number line? \emph{(asked)}
        \item \textbf{5.1} --- radicals are exponents in disguise
        \item \textbf{5.2--5.3} --- simplify, operate, rationalize
        \item \textbf{5.4} --- graph and transform; the endpoint is the anchor
        \item \textbf{5.5} --- solve, and always check for extraneous roots
        \item \textbf{5.6} --- composition; order is part of the answer
        \item \textbf{5.7} --- inverses. \emph{(answered)}
      \end{itemize}
    \end{column}
    \begin{column}{0.44\textwidth}
      \footnotesize
      \begin{block}{Next: the Unit 5 test}
        Convert, simplify, rationalize, read a graph, transform, solve and check, compose and find a
        domain, build and \emph{justify} an inverse. \\[4pt]
        The practice test is in your packet.
      \end{block}
      \vspace{4pt}
      \begin{center}
        \textbf{$x^{2}$ has two arms. $x^{3}$ has one.} \\[3pt]
        Everything else followed.
      \end{center}
    \end{column}
  \end{columns}
\end{frame}

\end{document}
